\documentclass[10pt,a4paper]{article}
\usepackage{amsmath,amssymb,amsthm,graphicx}
\usepackage{hyperref,bbm,bm,tikz}
\usetikzlibrary{decorations.pathmorphing}
\newtheorem{theorem}{Theorem}
\newtheorem{corollary}[theorem]{Corollary}
\newtheorem{lemma}[theorem]{Lemma}
\newtheorem{observation}[theorem]{Observation}
\newtheorem{proposition}[theorem]{Proposition}
\newtheorem{definition}[theorem]{Definition}
\newtheorem{claim}[theorem]{Claim}
\newtheorem{fact}[theorem]{Fact}
\newtheorem{remark}[theorem]{Remark}
\newtheorem{question}[theorem]{Question}

% General Tricks
\newcommand{\prob}[2]{\mathop{\mathrm{Pr}}_{#1}[#2]}
\newcommand{\avg}[2]{\mathop{\textbf{E}}_{#1}[#2]}
\newcommand{\poly}{\mathop{\mathrm{poly}}}
\newcommand{\tm}[1]{\textrm{#1}}
\newcommand{\bb}[1]{\mathbb{#1}}
\newcommand{\bra}[1]{\{#1\}}
\newcommand{\abs}[1]{\left|#1\right|}
\newcommand{\setcond}[2]{\left\{#1\: \middle|\: #2\right\}}
\newcommand{\errnote}[1]{{ \small #1}} 

% Complexity shortcuts
% General Math shortcuts
\newcommand{\F}{\mathbb{F}}
\newcommand{\R}{\mathbb{R}}
\newcommand{\C}{\mathbb{C}}
\newcommand{\AC}{\mathrm{AC}}
\newcommand{\naturals}{\mathbb{N}}
\newcommand{\integer}{\mathbb{Z}}
\newcommand{\rational}{\mathbb{Q}}
\newcommand{\ip}[2]{\langle #1, #2 \rangle}

\newcommand{\union}{\cup}
\newcommand{\intersect}{\cap}
\newcommand{\prefix}{\sqsupset}
\newcommand{\SPS}{\Sigma\Pi\Sigma}
\newcommand{\spsp}[2]{\Sigma\Pi^{#1}\Sigma\Pi^{#2}}
%\newcommand{\sPsP}[2]{\Sigma\Pi^{[#1]}\Sigma\Pi^{\{#2\}}}
\newcommand{\IMM}{\mathrm{IMM}}
\newcommand{\Det}{\mathrm{Det}}
\newcommand{\pder}[2]{\ensuremath{\partial{#1} \over
  \partial{#2}}}
\newcommand{\mc}[1]{\mathcal{#1}}
\newcommand{\Dim}{\mathrm{dim}}
\newcommand{\spd}[3]{\langle \partial_{#2} #1 \rangle_{\leq #3}}
\newcommand{\pd}[2]{\partial_{#2} #1}
\newcommand{\LM}{\mathrm{LM}}
\newcommand{\MSupp}{\mathrm{MSupp}}
\newcommand{\mynorm}[1]{\Vert #1 \Vert}

\newenvironment{myq}{\newline\noindent\begin{minipage}{\textwidth}}{\end{minipage}}


%margin comments
\newcommand{\comment}[2]{\marginpar{\tiny{\textbf{#1: }\textit{#2}}}}
\newcommand{\srikanth}[1]{\comment{S}{#1}}
\newcommand{\herve}[1]{\comment{H}{#1}}


\newcommand{\pdeg}{\mathrm{pdeg}}
\newcommand{\spec}{\mathrm{Spec}\,}
\newcommand{\per}{\mathrm{per}}
\newcommand{\sB}{s\mathcal{B}}
\newcommand{\B}{\mathcal{B}}
\newcommand{\Maj}{\mathrm{Maj}}
\newcommand{\Thr}{\mathrm{Thr}}
\newcommand{\MOD}{\mathrm{MOD}}
\newcommand{\OR}{\mathrm{OR}}
\newcommand{\AND}{\mathrm{AND}}
\newcommand{\EThr}{\mathrm{EThr}}


\newcommand{\mbf}{\mathbf}
\newcommand{\tbf}{\textbf}
\newcommand{\mb}{\mathbb}
\newcommand{\tx}{\text}
\newcommand{\Or}{\tx{OR}}
\newcommand{\Ex}{\tx{EX}}
\newcommand{\Near}{\tx{NEAR}}
\newcommand{\ol}{\overline}
\newcommand{\modulo}[1]{\,(\text{mod }#1)}
\newcommand{\lcm}{\tx{lcm}}
\newcommand{\charac}{\tx{char}}
\newcommand{\sgn}{\tx{sgn}}
\newcommand{\rank}{\mathrm{rank}}
\newcommand{\relrk}{\mathrm{relrk}}
\newcommand{\bias}{\mathrm{bias}}

\newcommand{\mP}{\mathcal{P}}
\newcommand{\mN}{\mathcal{N}}

% Shortcut usefull for the homogeneization step
\newcommand{\WE}{\mathit{WESym}}
\newcommand{\WP}{\mathit{WPow}}
\newcommand{\WQ}{\mathit{WQ}}
\newcommand{\setA}{\mathcal{A}}
\newcommand{\setB}{\mathcal{B}}
\newcommand{\setD}{\mathcal{D}}
\newcommand{\setZ}{\mathcal{Z}}

%Comments
%\newcommand{\comment}[2]{\marginpar{\tiny{\textbf{#1: }\textit{#2}}}}
\newcommand{\gm}[1]{\comment{GM}{#1}}
\newcommand{\hf}[1]{\comment{HF}{#1}}
\newcommand{\nl}[1]{\comment{NL}{#1}}
\newcommand{\st}[1]{\comment{ST}{#1}}
\newcommand{\ssr}[1]{\comment{SS}{#1}}

%Comments which do not fit in the margin
\newcommand{\commentL}[2]{\marginpar{\tiny{\textbf{#1}}} \par{\tiny{\textit{#2}}}}
\newcommand{\hfL}[1]{\commentL{HF}{#1}}
\title{Optimal logarithmic-degree depth for monotone algebraic formulas}
\author{Hervé Fournier, Nutan Limaye, Guillaume Malod, Srikanth Srinivasan, Sébastien Tavenas}
\date{doi:10.4230/LIPIcs.CCC.2023.28; CC BY4.0}
\begin{document}
\maketitle
\noindent\textit{Editorial note.} Original authors: Hervé Fournier, Nutan Limaye, Guillaume Malod, Srikanth Srinivasan, Sébastien Tavenas. arXiv:2302.06984v1 (2023-02-14), https://arxiv.org/abs/2302.06984v1. Creative Commons Attribution4.0, https://creativecommons.org/licenses/by/4.0/.\par\medskip
\noindent\textit{Editorial note.} Selected complete author native main statements/proofs, necessary context, original TikZ code and complete bibliography in standalone layout. Unused frontmatter/packages omitted; no mathematical changes or new proofs. Original source is CC BY4.0.\par\medskip
\noindent\textit{Allowed standard prerequisites (editorial).}\begin{itemize}
\item Standard formula depth reduction and Raz homogenization. Exact prior theorem statements in author preliminaries retained; near-linear Bonet-Buss proof also retained where used.
\end{itemize}\medskip
\begin{theorem}[Lower Bound]
\label{thm:lowerbound-intro}
Let $n$ and $d = d(n)$ be growing parameters such that $d(n)\leq \sqrt{n}$. Then there is a monotone algebraic formula $F$ of size at most $n$ and depth $O(\log d)$ computing a polynomial $P\in \F[x_1,\ldots,x_n]$ of degree at most $d$ such that any monotone formula $F'$ of depth $o(\log d)$ computing $P$ must have size $n^{\omega(1)}.$
\end{theorem}

\section{Preliminaries}
\label{sec:prelims}

\paragraph{Basic notation.} Throughout, unless otherwise specified, we work with polynomials over a field $\F$. We will work with the multivariate ring of polynomials $\F[x_1,\ldots,x_n]$ or its \emph{non-commutative} analog $\F\langle  x_1,\ldots,x_n \rangle.$ 

\subsection{Algebraic formulas} 
\label{sec:formula-prelims}

We start with some brief definitions and results related to algebraic formulas. For much more about this model, see the standard references~\cite{SY,Ramprasadsurvey}.

\paragraph{The model.} An algebraic formula over the multivariate polynomial ring $\F[x_1,\ldots,x_n]$ is a rooted, directed tree with edges directed towards the root. Leaves are labelled by variables $x_1,\ldots,x_n$ or by the constant \(1\) and edges by non-zero field constants. Internal nodes (i.e.,~gates) by $+$ and $\times$ and compute linear combinations (based on the edge weights) or products of their children. We will assume, with loss of generality, that if a node \(\alpha\) has for child a leaf labelled by \(1\), then \(\alpha\)  is a \(+\)-gate and that if a \(+\)-gate \(\alpha\) has only children labelled by \(1\), then \(\alpha\) is the output of the formula.\footnote{This ensures that a formula can compute polynomials with a constant term but forbids using many arithmetic operations just to compute constants.}
%\gm{added a footnote to give some rationale for the slightly uncommon definition}
A \emph{non-commutative} algebraic formula  over the multivariate polynomial ring $\F\langle x_1,\ldots,x_n\rangle$ is defined similarly, with the additional assumption that the children of any $\times$-gate are linearly ordered, and the corresponding product is computed in this order.

Unless explicitly stated, the algebraic formulas we consider have \emph{unbounded} fan-in (i.e.,~a gate can have any number of inputs). The \emph{size} of $F$ will denote the number of leaves,\footnote{\label{ft:fanin}This is within a constant factor of the number of gates, as long as each gate has fan-in at least $2$ each (which is without loss of generality).} the \emph{depth} of $F$ the longest leaf-to-root path. The \emph{product-depth} and the \emph{sum-depth} of $F$ are defined to be the maximum number of product gates and sum gates encountered on a leaf-to-root path, respectively.

A \emph{parse tree} of a formula $F$ is a subformula of $F$ which corresponds to the way a monomial is built in the evaluation of $F$.
Parse trees of $F$ can be defined inductively as follows:
\begin{itemize}
\item If $F$ has a top $+$-gate,
a parse tree of $F$ is obtained by taking a parse tree of one of its children together with the corresponding edge to the root of $F$;
\item If $F$ has a top $\times$-gate,
a parse tree of $F$ is obtained by a taking a parse tree of each of its children, together with the incoming edges of the root of $F$;
\item The only parse tree of a leaf is itself.
\end{itemize}
The polynomial computed by a parse tree is a single monomial, which
is equal to the product of the variables labelling its leaves, multiplied
by the product of the scalars labelling its edges.
The polynomial computed by a formula $F$ is easily seen to be the sum of the monomials computed by all its parse trees.

A formula $F$ is called \emph{monotone} if any monomial computed by a parse tree of $F$ has a non-zero coefficient
in the polynomial computed by $F$. 


We now recall some well-known results from the literature regarding algebraic formulas. It should be noted that these results (specifically Theorems~\ref{thm:standardDR}, \ref{thm:Raz} and \ref{thm:BCE} later on) are usually proved in the general setting of commutative formulas. However, it is easy to see that the proofs of these results carry over to the monotone, non-commutative setting without significant change.

\paragraph{Depth-reduction.} Classical results~\cite{BKM,Brent} show that any algebraic formula of small size can be simulated by one of small depth and not much larger size. Formally, 


\begin{theorem}
\label{thm:standardDR}
 Let $F$ be a (non-commutative or commutative) algebraic formula of size $s$ computing a polynomial $P$. Then there is an algebraic formula $F'$ of size at most $\poly(s)$ and depth $\Delta = O(\log s)$ computing $P$. We may also assume that each gate in $F'$ has fan-in $2$. Furthermore, $F$ is homogeneous and/or monotone, then so is $F'$.
\end{theorem}


\paragraph{Homogeneity.} Each gate in an algebraic formula has a \emph{syntactic degree} defined in a natural way. Leaves labelled by the constant~$1$ have syntactic degree~\(0\),  leaves labelled with a variable have syntactic degree~$1$, $\times$-gates have a syntactic degree that is the sum of the syntactic degrees of their children, and $+$-gates have a syntactic degree that is equal to the largest of the syntactic degrees of their children.
 The syntactic degree of a formula is defined as the syntactic degree of its output. Notice that in a formula  the syntactic degree of any gate is bounded by the syntactic degree of the formula. 

We will further assume that no gate computes the zero polynomial.

A formula is \emph{homogeneous} if each gate in the formula computes a homogeneous polynomial. Equivalently, in terms of syntactic degrees, this means that all the children of a sum gate have the same syntactic degree.

Raz~\cite{Raztensor} showed how to convert a possibly inhomogeneous formula $F$ to a homogeneous formula with a relatively small blow-up in size.

\begin{theorem}
\label{thm:Raz}
 Let $F$ be a (non-commutative or commutative) algebraic formula of size $s$ and product-depth $\Delta$ computing a polynomial $P$ such that all gates in $F$ have fan-in $2$. Then there is a \emph{homogeneous} algebraic formula $F'$ of size at most $O\left(s\cdot \binom{\Delta + d+1}{d}\right)$ and product-depth $\Delta$ computing $P$.  In particular, if $\Delta = O(\log s)$ and $d = O(\log s)$, then the formula $F'$ has size $\poly(s).$
\end{theorem}



\begin{theorem}[Depth-reduction with near-linear size]
	\label{thm:BCE}
	The following holds for any $\varepsilon > 0.$ Let $F$ be a (non-commutative or commutative) algebraic formula of size $s$ computing a polynomial $P$. Then there is an algebraic formula $F'$ of size at most $s^{1+\varepsilon}$ and depth $\Delta = 2^{O(1/\varepsilon)}\cdot \log s$ computing $P$. We may also assume that each gate in $F'$ has fan-in $2$. Furthermore, if $F$ is homogeneous and/or monotone, then so is $F'$.
\end{theorem}

\begin{theorem}
	\label{thm:lowerbound}
	Let $n$ and $d = d(n)$ be growing parameters such that $d(n)\leq \sqrt{n}$. Then there is a monotone algebraic formula $F$ of size at most $n$ and depth $O(\log d)$ computing a polynomial $P\in \F[x_1,\ldots,x_n]$ of degree at most $d$ such that any monotone formula $F'$ of depth $o(\log d)$ computing $P$ must have size $n^{\omega(1)}.$
\end{theorem}

\section{Tightness}

\label{sec:tightness}


Given integers $k \geq 1$ and $r \geq 2$, we will define a polynomial $H^{(k,r)}$.
Intuitively, we want to define this polynomial as a standard universal polynomial for formulas. It is composed of \(k\)-nested inner products, each one of size \(r\).
In the following we will drop the superscript in $H^{(k,r)}$ and write simply $H$ instead.

The polynomial $H$ will be defined over the set of
%\hf{added the number of variables}
$(2r)^k$  variables
$$\{x_{\sigma,\tau}\ |\ \sigma\in [2]^k, \tau \in [r]^{k}\}.$$
Let us define recursively polynomials $H_{u,v}$ for all $(u,v) \in [2]^{\leq k} \times [r]^{\leq k}$ such that $|u|=|v|$:
%\hf{Changed to $H_{u,v} = x_{u,v}$ when $|u|=|v|=k$. Can be reversed.}
\begin{align*}
	H_{u,v} = & \ x_{u,v} 	&& \text{when }|u|=|v|=k\\
	H_{u,v} = & \sum_{a=1}^{r} H_{u 1,v a} H_{u 2, v a} && \text{otherwise}.
\end{align*}
The polynomial $H$ is defined as the polynomial \(H_{\varepsilon,\varepsilon}\).
Note that $H$ is a polynomial of degree $d=2^k$ and has $r^{d-1}$ monomials.

From its definition, $H$ is computed by a monotone formula $M$
%\hf{added the size of $M$}
of size $(2r)^k$
and depth $2k$,
with a $+$-gate at the top, alternating layers of $+$-gates and $\times$-gates, with $+$-gates of fan-in $r$
and $\times$-gates of fan-in $2$, and leaves labelled with distinct variables.

For words $u$ and $v$ over the same alphabet, we write $u \prefix v$ if $u$ is a prefix of $v$.
There is a natural one-to-one correspondance between prefixes of words of $([r]\times[2])^k$
and nodes of $M$, which is the following.
Let $\sigma = \sigma_1 \ldots \sigma_k \in [2]^k$, and $\tau = \tau_1 \ldots \tau_{k} \in [r]^{k}$.
The word $\tau_1 \sigma_1 \ldots \tau_k \sigma_k$ corresponds to a path from the root of $M$
to the leaf labelled $x_{\sigma,\tau}$, while proper prefixes of $\tau_1 \sigma_1 \ldots \tau_k \sigma_k$
correspond to internal gates in $M$ along this path.
For $\ell < k$, $u=u_1\ldots u_\ell \in [2]^\ell$ and $v=v_1 \ldots v_\ell \in [r]^\ell$,
the node which corresponds to the word $v_1 u_1 \ldots v_\ell u_\ell$ is the $+$-gate of $M$ computing $H_{u,v}$.
%\hf{Thought about a figure I could do there if needed (Caramels, bonbons et chocolats...).}

The polynomial $H$ is easily seen to be set-multilinear with respect to
the sets of variables $\{ X_\sigma\ |\ \sigma \in [2]^k\}$
where $X_\sigma = \{x_{\sigma,\tau}\ |\ \tau \in [r]^{k} \}$. This means that each monomial has exactly one variable from each set $X_\sigma$.

\begin{remark}
	\label{rem:isomorphic-subpolynomials}
	For $|u|=|v|=\ell \leq k$, $u \in [2]^\ell$ and $v \in [r]^\ell$,
	the polynomial $H^{(k,r)}_{u,v}$
	is defined over the set of variables
	$$X_{u,v} = \{ x_{\sigma,\tau}\ |\ \sigma \in [2]^k,\ \tau \in [r]^{k}, u \prefix \sigma,\ v\prefix \tau\}$$
	and is the polynomial $H^{(k-\ell,r)}$ (upto renaming of the variables).
	In particular, its degree is $d'=2^{k-\ell}$,
	it has $r^{d'-1}$ monomials 
	and it is set-multilinear with
	respect to $\{ X_\sigma\ |\ \sigma \in [2]^k,\ u\prefix \sigma\}$.
\end{remark}

Before proving hardness of $H$ for small-depth monotone formulas, we need to show that the gates of a monotone formula computing~\(H\) cannot compute too many monomials. This is proved in Lemma~\ref{lem:number-of-monomials} below.

\begin{proposition}
	\label{prop:product-of-variables-in-H}
	Consider two variables $x_{\sigma,\tau}$ and $x_{\sigma',\tau'}$.
	If $x_{\sigma,\tau} x_{\sigma',\tau'}$ appears in a monomial of~$H$,
	and if $\sigma$ and $\sigma'$ have a common prefix of length $\ell < k$,
	then $\tau$ and $\tau'$ have a common prefix of length $\ell+1$.
\end{proposition}

\begin{proof}
	Observe that
	$x_{\sigma_1,\tau_1} \ldots x_{\sigma_p,\tau_p}$ is a monomial of $H$ if and only if
	these variables
	form the leaves of a parse tree of $M$.
	As observed above, the root-to-leaf path leading to the variable
	$x_{\sigma_i,\tau_i}$ in $M$
	is obtained by taking in turn the first letter of  $\tau_i$, the first letter of $\sigma_i$, the second letter of $\tau_i$, etc.
	
	If the product $x_{\sigma,\tau} x_{\sigma',\tau'}$ appears in a monomial of $H$, it must
	be possible to complete
	the union of the two paths, from the root to $x_{\sigma,\tau}$ and from the root to $x_{\sigma',\tau'}$,
	into a parse tree of the formula $M$.
	
	If the longest common prefix of $\tau$ and $\tau'$ were of length at most $\ell$,
	%\hf{``of length at most $\ell<k$'' $\rightarrow$ ``of length $\ell<k$''.}\st{$\ell$ is fixed by the statement of the proposition, so I would prefer this.}
	then the lowest common ancestor of $x_{\sigma,\tau}$ and $x_{\sigma',\tau'}$ in $M$ would be a $+$-gate,
	which is not possible in a parse tree where each $+$-gate has a single child.
\end{proof}

\begin{lemma}\label{lem:number-of-monomials}
	If \(F\) is a monotone formula which computes the polynomial \(H\) and if \(\alpha\) is a gate of \(F\) of degree \(d_{\alpha}\), then the number of monomials of the polynomial computed at gate \(\alpha\) is at most \(r^{d_{\alpha}-1}\).
\end{lemma}

\begin{proof}
	Since \(H\) has \(r^{d-1}\) monomials, the result is true when \(d_{\alpha} = d\) by monotonicity. Assume now that \(d_{\alpha} < d\). 
	
	Let $$I = \{ \sigma \in [2]^k\ |\ \text{some variable of $X_\sigma$ appears in $\alpha$}\}.$$
	For $u \in [2]^{\leqslant k}$, let $I_u = \{ \sigma \in [2]^k\ |\ u \prefix \sigma\}$.
	Let $\{u_1, \ldots, u_p\}$ be the set of words $w$ of minimal length such that $I_w \subseteq I$.
	Then $I$ is the disjoint union $\bigcup_{\ell \in [p]} I_{u_\ell}$.
	Since $d_\alpha<d$, $I \neq [2]^k$ so no $u_\ell$ is the empty word.
	
	Consider some $\ell \in [p]$.
	Let $\bar{u}_\ell$ be obtained by switching the last letter of $u_\ell$.
	%By maximality $I_{\bar{u}_\ell} \nsubseteq I$: 
	By minimality of the length of $u_\ell$, we must have $I_{\bar{u}_\ell} \nsubseteq I$. 
	Let $\sigma$ be a word in   $I_{\bar{u}_\ell} \setminus I$. 
	Since $F$ is monotone and computes a set-multilinear polynomial, it is a set-multilinear formula
	and therefore the polynomial computed at $\alpha$
	must be multiplied by some variable $x_{\sigma,\tau}$.
	Let $v_\ell$ be the prefix of length $|u_\ell|$ of $\tau$.
	Consider any variable $x_{\sigma',\tau'}$ appearing in $\alpha$ such that $u_\ell \prefix \sigma'$.
	Since the product $x_{\sigma,\tau} x_{\sigma',\tau'}$ must appear in a monomial of $H$ by monotonicity,
	it must be that $v_\ell \prefix \tau'$
	by Proposition~\ref{prop:product-of-variables-in-H}.
	
	
	Any monomial $m$ in the polynomial computed in gate $\alpha$ can be written in a unique way $m=m_1 \cdots m_p$ with $m_\ell$
	set-multilinear with respect to $\{X_{\sigma}\ |\ u_\ell \prefix \sigma\}$.
	By the above, $m_\ell$ is a monomial over the variables $X_{u_\ell,v_\ell}$ of degree $|I_{u_\ell}|$.
	By monotonicity, it should be possible to complete leaves of $M$ labelled with variables from $m_\ell$
	into a parse tree of $M$ appearing in $H$, which proves that 
	$m_\ell$ is a monomial of $H_{u_\ell,v_\ell}$.
	There are at most $r^{|I_{u_\ell}|-1}$ such submonomials $m_\ell$ by Remark~\ref{rem:isomorphic-subpolynomials}.
	It follows that the number of monomials of the polynomial computed in node $\alpha$ is at most
	\[
	\prod_{\ell=1}^p r^{\lvert I_{u_{\ell}} \rvert -1} \leq r^{\lvert I \rvert -1} = r^{d_\alpha-1}.\qedhere
	\]
\end{proof}	

We are ready to prove hardness of the polynomial $H$ for monotone computation.
We shall make use of the following ``product lemma'', which comes in different forms in e.g.~\cite{SY, HrubesY11,Ramprasadsurvey}.

\begin{lemma}\label{lem:product-decomposition}
	A degree-$d$ homogeneous formula $F$ of size $s$ and product-depth $\Delta$
	can be written as a sum of $O(s)$ polynomials, % $F_1, \ldots , F_{s'}$,
	each of which
	is a product of $\Omega(\Delta d^{1/\Delta})$ many polynomials of
	positive degree. Moreover, each of these polynomials is computed by
	some gate in $F$.
\end{lemma}

\begin{proposition}
	\label{prop:hardness-of-H}
	If \(F\) is a monotone formula of product-depth \(\Delta  \leq \log d\)
	which computes \(H\), then its size is at least $r^{\Omega(\Delta d^{1/\Delta})}$.
\end{proposition}

\begin{proof}
	Let $t=\Delta d^{1/\Delta}$.
	Since $F$ is monotone, it is homogeneous and by Lemma~\ref{lem:product-decomposition} can be written as
	$$F = \sum_{i=1}^{s'} \prod_{j=1}^{t_i} F_{i,j}$$
	with $s'=O(s)$ and $t_i =\Omega(t)$ for all $i$, 
	and $F_{i,j}$ is of degree at least $1$ and computed by some gate in $F$.
	
	Let $d_{i,j}$ be the degree of $F_{i,j}$.
	By Lemma~\ref{lem:number-of-monomials}, each $F_{i,j}$ computes at most $r^{d_{i,j}-1}$ monomials.
	The number of monomials of $\prod_{i=1}^{t_j} F_{i,j}$ is therefore bounded by
	$$\prod_{j=1}^{t_i} r^{d_{i,j}-1} \leq r^{\sum_{j=1}^{t_i} d_{i,j}-t_i} \leq r^{d-t}$$
	since $\sum_{j=1}^{t_i} d_{i,j} = d$.
	It follows that the number of monomials computed by $F$ is at most $s \cdot r^{d-t}$.
	Since $F$ computes $H$ which has $r^{d-1}$ monomials, we get $s = r^{\Omega (t)}$.
\end{proof}

We can now get Theorem~\ref{thm:lowerbound-intro} from the introduction, which is restated here for convenience.

\begin{proof}[Proof of Theorem~\ref{thm:lowerbound-intro}]
	Choose parameters $k(n)$ and $r(n)$ 
	such that $P(n):=H^{(k,r)}$ has $\Theta(n)$ variables and degree $\Theta(d)$: let
	$k=\log d$ and $r=\frac{1}{2} n^{1/\log d}$.
	Condition $d(n)\leq \sqrt{n}$ ensures that $r\geqslant 2$.
	The polynomial $P$ has a monotone formula of size $O(n)$ and depth $O(\log d)$.
	By Proposition~\ref{prop:hardness-of-H}, 
	any monotone formula of product-depth $\Delta \leq \log d$ computing $P$
	has size
	$$r^{\Omega(\Delta d^{1/\Delta})} = 
	\left(\frac{1}{2} n^{1/\log d}\right)^{\Omega(\Delta d^{1/\Delta})} 
	= (n/d)^{\left( \frac{\Delta d^{1/\Delta}}{\log d} \right)}$$
	which is $n^{\Omega \left( \frac{\Delta d^{1/\Delta}}{\log d} \right)}$
	using the hypothesis $d(n)\leq \sqrt{n}$.
	Since $\frac{\Delta d^{1/\Delta}}{\log d} \rightarrow +\infty$ when $\Delta = o \left(\log d \right)$ this bound is $n^{\omega(1)}.$
\end{proof}

\begin{thebibliography}{GKKS16}

\bibitem[AV08]{AV}
Manindra Agrawal and V.~Vinay.
\newblock Arithmetic circuits: A chasm at depth four.
\newblock In {\em proceedings of Foundations of Computer Science (FOCS)}, pages
  67--75, 2008.

\bibitem[B\"00]{BurgisserVPVNP}
Peter B\"{u}rgisser.
\newblock Cook's versus {V}aliant's hypothesis.
\newblock {\em Theoretical Computer Science}, 235(1):71--88, 2000.

\bibitem[BB94]{BonetBuss}
Maria~Luisa Bonet and Samuel~R. Buss.
\newblock Size-depth tradeoffs for boolean formulae.
\newblock {\em Information Processing Letters}, 49(3):151 -- 155, 1994.

\bibitem[BCE95]{BCE}
Nader~H. Bshouty, Richard Cleve, and Wayne Eberly.
\newblock Size-depth tradeoffs for algebraic formulas.
\newblock {\em {SIAM} J. Comput.}, 24(4):682--705, 1995.

\bibitem[BKM73]{BKM}
R.~Brent, D.~Kuck, and K.~Maruyama.
\newblock The parallel evaluation of arithmetic expressions without division.
\newblock {\em IEEE Transactions on Computers}, C-22(5):532--534, 1973.

\bibitem[Bre74]{Brent}
Richard~P. Brent.
\newblock The parallel evaluation of general arithmetic expressions.
\newblock {\em Journal of the ACM}, 21(2):201--206, April 1974.

\bibitem[CKSV16]{CKSV}
Suryajith Chillara, Mrinal Kumar, Ramprasad Saptharishi, and V.~Vinay.
\newblock The chasm at depth four, and tensor rank : Old results, new insights.
\newblock {\em CoRR}, abs/1606.04200, 2016.

\bibitem[CLS18]{CLS}
Suryajith Chillara, Nutan Limaye, and Srikanth Srinivasan.
\newblock Small-depth multilinear formula lower bounds for iterated matrix
  multiplication, with applications.
\newblock In {\em {STACS}}, volume~96 of {\em LIPIcs}, pages 21:1--21:15.
  Schloss Dagstuhl - Leibniz-Zentrum fuer Informatik, 2018.

\bibitem[GKKS16]{GKKSdepth3}
Ankit Gupta, Pritish Kamath, Neeraj Kayal, and Ramprasad Saptharishi.
\newblock Arithmetic circuits: {A} chasm at depth 3.
\newblock {\em {SIAM} Journal of Computing}, 45(3):1064--1079, 2016.

\bibitem[HY11]{HrubesY11}
Pavel Hrube\v{s} and Amir Yehudayoff.
\newblock Homogeneous formulas and symmetric polynomials.
\newblock {\em Comput. Complexity}, 20(3):559--578, 2011.

\bibitem[KdOS19]{KOS19}
Mrinal Kumar, Rafael~Mendes de~Oliveira, and Ramprasad Saptharishi.
\newblock Towards optimal depth reductions for syntactically multilinear
  circuits.
\newblock In Christel Baier, Ioannis Chatzigiannakis, Paola Flocchini, and
  Stefano Leonardi, editors, {\em 46th International Colloquium on Automata,
  Languages, and Programming, {ICALP} 2019, July 9-12, 2019, Patras, Greece},
  volume 132 of {\em LIPIcs}, pages 78:1--78:15. Schloss Dagstuhl -
  Leibniz-Zentrum f{\"{u}}r Informatik, 2019.

\bibitem[Koi12]{Koiran}
Pascal Koiran.
\newblock Arithmetic circuits: The chasm at depth four gets wider.
\newblock {\em Theor. Comput. Sci.}, 448:56--65, 2012.

\bibitem[Kos86]{Kos86}
S~Rao Kosaraju.
\newblock Parallel evaluation of division-free arithmetic equations.
\newblock In {\em Proceedings of the Eighteenth Annual ACM Symposium on Theory
  of Computing}, pages 231--239, 1986.

\bibitem[KSS14]{KSS}
Neeraj Kayal, Chandan Saha, and Ramprasad Saptharishi.
\newblock A super-polynomial lower bound for regular arithmetic formulas.
\newblock In David~B. Shmoys, editor, {\em Symposium on Theory of Computing,
  {STOC} 2014, New York, NY, USA, May 31 - June 03, 2014}, pages 146--153.
  {ACM}, 2014.

\bibitem[LST21a]{LST-ECCC}
Nutan Limaye, Srikanth Srinivasan, and S{\'{e}}bastien Tavenas.
\newblock Superpolynomial lower bounds against low-depth algebraic circuits.
\newblock {\em Electron. Colloquium Comput. Complex.}, {TR21-081}, 2021.

\bibitem[LST21b]{LST-FOCS}
Nutan Limaye, Srikanth Srinivasan, and S{\'{e}}bastien Tavenas.
\newblock Superpolynomial lower bounds against low-depth algebraic circuits.
\newblock In {\em 62nd {IEEE} Annual Symposium on Foundations of Computer
  Science, {FOCS} 2021, Denver, CO, USA, February 7-10, 2022}, pages 804--814.
  {IEEE}, 2021.

\bibitem[Nis91]{Nisan}
Noam Nisan.
\newblock Lower bounds for non-commutative computation.
\newblock In {\em Proceedings of the twenty-third annual ACM symposium on
  Theory of computing}, pages 410--418, 1991.

\bibitem[NW97]{NW95}
Noam Nisan and Avi Wigderson.
\newblock Lower bounds on arithmetic circuits via partial derivatives.
\newblock {\em Computational Complexity}, 6(3):217--234, 1997.

\bibitem[PM75]{PreparataMuller1}
Franco~P. Preparata and David~E. Muller.
\newblock The time required to evaluate division-free arithmetic expressions.
\newblock {\em Inf. Process. Lett.}, 3(5):144--146, 1975.

\bibitem[PM76]{PreparataMuller2}
Franco~P. Preparata and David~E. Muller.
\newblock Efficient parallel evaluation of boolean expression.
\newblock {\em {IEEE} Trans. Computers}, 25(5):548--549, 1976.

\bibitem[Raz06]{Raznc2nc1}
Ran Raz.
\newblock Separation of multilinear circuit and formula size.
\newblock {\em Theory of Computing}, 2(1):121--135, 2006.

\bibitem[Raz13]{Raztensor}
Ran Raz.
\newblock Tensor-rank and lower bounds for arithmetic formulas.
\newblock {\em Journal of the ACM}, 60(6):40:1--40:15, 2013.

\bibitem[RY08]{RYbalance}
Ran Raz and Amir Yehudayoff.
\newblock Balancing syntactically multilinear arithmetic circuits.
\newblock {\em Computational Complexity}, 17(4):515--535, 2008.

\bibitem[Sap15]{Ramprasadsurvey}
Ramprasad Saptharishi.
\newblock A survey of lower bounds in arithmetic circuit complexity.
\newblock Github survey, 2015.

\bibitem[Spi71]{Spira71tradeoffs}
Philip~M. Spira.
\newblock On time hardware complexity tradeoffs for {B}oolean functions.
\newblock In Shu Lin, editor, {\em Proceedings of the Fourth Hawaii
  International Conference on System Sciences}, pages 525--527. Western
  Periodicals Company, North Hollywood, California, 1971.

\bibitem[SS80]{ShamirSnir}
Eli Shamir and Marc Snir.
\newblock On the depth complexity of formulas.
\newblock {\em Math. Syst. Theory}, 13:301--322, 1980.

\bibitem[SW01]{SW}
Amir Shpilka and Avi Wigderson.
\newblock Depth-3 arithmetic circuits over fields of characteristic zero.
\newblock {\em Computational Complexity}, 10(1):1--27, 2001.

\bibitem[SY10]{SY}
Amir Shpilka and Amir Yehudayoff.
\newblock Arithmetic circuits: A survey of recent results and open questions.
\newblock {\em Foundations and Trends in Theoretical Computer Science},
  5:207--388, March 2010.

\bibitem[Tav15]{Tav13}
S{\'{e}}bastien Tavenas.
\newblock Improved bounds for reduction to depth 4 and depth 3.
\newblock {\em Information and Computation}, 240:2--11, 2015.

\bibitem[TLS22]{LST-STOC}
S{\'{e}}bastien Tavenas, Nutan Limaye, and Srikanth Srinivasan.
\newblock Set-multilinear and non-commutative formula lower bounds for iterated
  matrix multiplication.
\newblock In {\em {STOC}}, pages 416--425. {ACM}, 2022.

\bibitem[VSBR83]{VSBR}
Leslie~G. Valiant, Sven Skyum, Stuart~J. Berkowitz, and Charles Rackoff.
\newblock {Fast Parallel Computation of Polynomials Using Few Processors}.
\newblock {\em SIAM Journal of Computing}, 12(4):641--644, 1983.

\bibitem[YK68]{YK68}
SV~Yablonskii and VP~Kozyrev.
\newblock Mathematical problems of cybernetics.
\newblock {\em Information Materials of Scientific Council of Akad. Nauk SSSR
  on Complex Problem “Kibernetika}, 19:3--15, 1968.

\end{thebibliography}
\end{document}
